\SourcePage{746}{749}
\section{Inelastic Collisions of Fast Electrons with Atoms}

Inelastic collisions of fast electrons with atoms can be treated in the Born
approximation in the same way as the elastic collisions considered in
\S~139\footnote{Most of the results presented in \S\S~148--150 were obtained
by Bethe (H.~A.~Bethe, 1930).}. The applicability condition for the Born
approximation still requires the incident electron's speed to be large in
comparison with the speeds of the atomic electrons. The energy lost in the
collision, however, may have any value. If the electron loses a substantial
part of its energy, the atom is ionized and the energy is transferred to one
of its electrons. Yet we may always designate as the scattered electron the
one of the two that has the higher speed after the collision. Thus, when the
incident electron is fast, the scattered electron will also be fast.

For collisions between an electron and an atom, the frame in which their
centre of mass is at rest may be regarded,

\SourcePage{747}{750}
as already noted, as coinciding with the frame in which the atom is at rest;
it is this latter frame that will be used below.

An inelastic collision changes the internal state of the atom. The atom may
pass from its ground state to an excited state in either the discrete or the
continuous spectrum; in the latter case the atom is ionized. These cases can
be treated together in deriving the general formulae.

We start, as in \S~126, from the general formula for the transition
probability between continuum states and apply it to the system formed by the
incident electron and the atom. Let $\mathbf p$, $\mathbf p'$ be the momenta
of the incident electron, and let $E_0$, $E_n$ be the atomic energies before
and after the collision, respectively. In place of (126.9), the transition
probability is
\begin{equation}
 dw_n=\frac{2\pi}{\hbar}|\langle n,\mathbf p'|U|0,\mathbf p\rangle|^2
 \delta\!\left(\frac{p'^2-p^2}{2m}+E_n-E_0\right)
 \frac{d^3p'}{(2\pi\hbar)^3},
 \tag{148.1}
\end{equation}
where the matrix element is taken of the interaction energy between the
incident electron and the atom,
\[
 U=\frac{Ze^2}{r}-\sum_{a=1}^{Z}\frac{e^2}{|\mathbf r-\mathbf r_a|}
\]
($\mathbf r$ is the radius vector of the incident electron, $\mathbf r_a$
those of the atomic electrons, and the origin is at the atomic nucleus; $m$
is the electron mass).

The electron wave functions $\psi_{\mathbf p}$ and $\psi_{\mathbf p'}$ are
given by the same formulae (126.10), (126.11); $dw$ is then the collision
cross-section $d\sigma$. Denote the atomic wave functions in the initial and
final states by $\psi_0$ and $\psi_n$. If the final atomic state belongs to
the discrete spectrum, $\psi_n$, like $\psi_0$, has the usual normalization
to unity. If instead the atom passes to a continuum state, the wave function
is normalized to a delta function of the parameters $\nu$ specifying such
states (these parameters may, for example, be the atomic energy and the
components of the momentum of the electron ejected from the atom upon
ionization). The resulting cross-sections give the probability of a collision
in which the atom passes to continuum states whose parameter values lie
between $\nu$ and $\nu+d\nu$.

Integration of (148.1) over the magnitude of $p'$ gives
\[
 d\sigma_n=\frac{mp'}{4\pi^2\hbar^4}|\langle n\mathbf p'|U|0\mathbf p\rangle|^2d\Omega',
\]

\SourcePage{748}{751}
where $p'$ is fixed by energy conservation:
\begin{equation}
 \frac{p^2-p'^2}{2m}=E_n-E_0.
 \tag{148.2}
\end{equation}
Substitution of the electron wave functions (126.10), (126.11) into the
matrix element gives
\begin{equation}
 d\sigma_n=\frac{m^2}{4\pi^2\hbar^4}\frac{p'}p
 \left|\iint Ue^{-i\mathbf q\mathbf r}\psi_n^*\psi_0\,d\tau\,dV\right|^2d\Omega
 \tag{148.3}
\end{equation}
($d\tau=dV_1dV_2\ldots dV_Z$ is the configuration-space element for the
$Z$ atomic electrons; the prime on $d\Omega$ has been omitted)\footnote{In
this form it is the general perturbation-theory formula, applicable not only
to collisions of electrons with atoms but to any inelastic collision of two
particles. It gives the scattering cross-section in the frame in which the
particles' centre of mass is at rest (in that case $m$ is the reduced mass of
the two particles).}. When $n=0$ and $p=p'$, formula (148.3) reduces to the
formula for the elastic-scattering cross-section.

Because $\psi_n$ and $\psi_0$ are orthogonal, the term in $U$ describing the
interaction $Ze^2/r$ with the nucleus vanishes upon integration over
$d\tau$. Thus, for inelastic collisions,
\begin{equation}
 d\sigma_n=\frac{m^2}{4\pi^2\hbar^4}\frac{p'}p
 \left|\sum_a\iint\frac{e^2}{|\mathbf r-\mathbf r_a|}e^{-i\mathbf q\mathbf r}
 \psi_n^*\psi_0\,d\tau\,dV\right|^2d\Omega.
 \tag{148.4}
\end{equation}

The integration over $dV$ can be performed as in \S~139. The integral
\[
 \varphi_{\mathbf q}(\mathbf r_a)=\int\frac{e^{-i\mathbf q\mathbf r}}{|\mathbf r-\mathbf r_a|}\,dV
\]
is formally the Fourier component of the potential produced at $\mathbf r$
by charges distributed in space with density
$\rho=\delta(\mathbf r-\mathbf r_a)$. Formula (139.1) therefore gives
\begin{equation}
 \varphi_{\mathbf q}(\mathbf r_a)=\frac{4\pi}{q^2}e^{-i\mathbf q\mathbf r_a}.
 \tag{148.5}
\end{equation}
Substitution in (148.4) yields the following final \emph{general} expression
for the cross-section of inelastic collisions:
\begin{equation}
 d\sigma_n=\left(\frac{e^2m}{\hbar^2}\right)^2\frac{4\kappa'}{\kappa q^4}
 \left|\left\langle n\left|\sum_a e^{-i\mathbf q\mathbf r_a}\right|0\right\rangle\right|^2d\Omega,
 \tag{148.6}
\end{equation}
where the matrix element is taken between atomic wave functions, and the
momenta have been replaced by the wave vectors
$\boldsymbol{\kappa}=\mathbf p/\hbar$,

\SourcePage{749}{752}
$\boldsymbol{\kappa}'=\mathbf p'/\hbar$. This formula gives the probability
of a collision in which the electron is scattered into the solid-angle
element $d\Omega$ and the atom passes to its $n$th excited state. The vector
$-\hbar\mathbf q$ is the momentum transferred by the electron to the atom in
the collision.

In calculations it is sometimes more convenient to refer the cross-section
to the element $dq$ of the magnitude of $\mathbf q$, rather than to a
solid-angle element. The vector $\mathbf q$ is defined by
$\mathbf q=\boldsymbol{\kappa}'-\boldsymbol{\kappa}$; its magnitude satisfies
\begin{equation}
 q^2=\kappa^2+\kappa'^2-2\kappa\kappa'\cos\vartheta.
 \tag{148.7}
\end{equation}
Hence, at fixed $\kappa$, $\kappa'$, that is, for a specified energy loss by
the electron,
\begin{equation}
 q\,dq=\kappa\kappa'\sin\vartheta\,d\vartheta
 =\frac{\kappa\kappa'}{2\pi}d\Omega.
 \tag{148.8}
\end{equation}
Formula (148.6) may therefore be rewritten as
\begin{equation}
 d\sigma_n=8\pi\left(\frac{e^2}{\hbar v}\right)^2\frac{dq}{q^3}
 \left|\left\langle n\left|\sum_a e^{-i\mathbf q\mathbf r_a}\right|0\right\rangle\right|^2.
 \tag{148.9}
\end{equation}

The vector $\mathbf q$ plays an essential part in the calculations below.
Let us examine more closely how it is related to the scattering angle
$\vartheta$ and to the energy $E_n-E_0$ transferred in the collision. We
shall see below that the principal contribution comes from collisions which
produce scattering through small angles ($\vartheta\ll1$) and transfer an
energy small compared with the incident electron's energy $E=mv^2/2$:
$E_n-E_0\ll E$. The difference $\kappa-\kappa'$ is then also small
($\kappa-\kappa'\ll\kappa$), and consequently
\[
 E_n-E_0=\frac{\hbar^2}{2m}(\kappa^2-\kappa'^2)
 \approx\frac{\hbar^2}{m}\kappa(\kappa-\kappa')
 =\hbar v(\kappa-\kappa').
\]
For small $\vartheta$, (148.7) gives
$q^2\approx(\kappa-\kappa')^2+(\kappa\vartheta)^2$, and hence
\begin{equation}
 q=\sqrt{\left(\frac{E_n-E_0}{\hbar v}\right)^2+(\kappa\vartheta)^2}.
 \tag{148.10}
\end{equation}
The minimum value of $q$ is
\begin{equation}
 q_{\min}=\frac{E_n-E_0}{\hbar v}.
 \tag{148.11}
\end{equation}

At small angles, one may further distinguish regions according to the ratio
of the small quantities $\vartheta$ and $v_0/v$ ($v_0$ is of the order of an
atomic-electron speed). If energy transfers of the order of the
atomic-electron energy $\varepsilon_0$ are considered
($E_n-E_0\sim\varepsilon_0\sim mv_0^2$), then for
$(v_0/v)^2\ll\vartheta\ll1$,
\begin{equation}
 q=\kappa\vartheta=mv\vartheta/\hbar.
 \tag{148.12}
\end{equation}

